\section{Conclusion}
\label{sec:conclusion}

The strongest evidence supports a conditional process--structure--conductivity relation in LLZO. Dopants and lithium inventory alter phase stability and carrier occupancy; precursor route and green-body history set the reaction length scale; thermal treatment, cover and crucible determine lithium retention, density and grain-boundary chemistry; and these states control how bulk, boundary and total resistance appear in impedance. Rapid sintering, membranes and interface layers expand the processing window, but their benefits remain conditional on measured phase fraction, relative density, surface chemistry, contact pressure and a named transport protocol \cite{anderson2024comprehensive,heywood2023tailoring,moy2025irreversible,yao2026elimination,zhang2024unveiling,klimpel2023standardizing}. A complete LLZO phase diagram and a universal route ranking are not established by the present evidence; the transferable result is the minimum comparison record.

The highest-value next experiments are therefore the three discriminating designs described above: first, map lithium-addition order, cover state and compact thickness with assay-corrected high-temperature diffraction; second, test Ta--LLZO grain-boundary chemistry at matched density with calibrated microscopy and a standardized critical-current protocol; third, quantify and remove the ultrafast-sintered surface layer using paired sister specimens. These experiments can be scheduled independently, and each has a clear failure criterion. Together they offer the fastest route to deciding whether a reported conductivity gain arises from phase stabilization, densification, boundary chemistry or simply a change in measurement conditions.
